\subsection{Neutrino oscillations and large extra dimensions --- arXiv:hep-ph/0002199}

\textbf{Authors:} Riccardo Barbieri, Paolo Creminelli, Alessandro Strumia

\paragraph{主な主張}
bulk右巻きニュートリノのKK塔が、通常の二状態振動とは異なる大気ニュートリノの消失パターンを生むことを示し、著者らの最小模型・超新星評価ではSN1987Aから$R\lesssim10^{-10}\,\mathrm{m}$程度の厳しい制約が生じると論じる\cite{Barbieri:2000BulkOscillations}。

\paragraph{新規性}
KK塔を含む真空・物質中の振動を弱混合に限定せず体系化し、多数の状態への転換がもたらす非標準的な減衰を明らかにした。

\paragraph{位置付け}
大気ニュートリノによるbulk振動の基礎文献であり、超新星に由来する数値制約は模型と輸送過程の仮定を伴うものとして読む必要がある。
